\chapter{Analyse, Conception et Choix Technologiques}

\noindent\textbf{Plan du chapitre}
\begin{enumerate}[nosep]
  \item Introduction
  \item Spécification des besoins
  \item Identification des acteurs
  \item Diagramme de cas d'utilisation global
  \item Architecture microservices
  \item Schéma de base de données
  \item Environnement de développement
  \item Choix technologiques justifiés
  \item Conclusion
\end{enumerate}

% ─────────────────────────────────────────────────────────────
\section{Introduction}
% ─────────────────────────────────────────────────────────────

Avant toute ligne de code, trois questions structurantes exigent une réponse
documentée~: \textit{quoi} construire (besoins fonctionnels et non fonctionnels)~;
\textit{comment} le décomposer (architecture)~; et \textit{avec quoi} (choix
technologiques comparés). Ce chapitre répond à ces trois questions en s'appuyant
sur des critères objectifs propres aux contraintes de GerAI~: budget zéro de
licences, déploiement on-premise, équipe réduite, calendrier de quatre mois.

% ─────────────────────────────────────────────────────────────
\section{Spécification des besoins}
% ─────────────────────────────────────────────────────────────

\subsection{Besoins fonctionnels~: Product Backlog}

\begin{table}[H]
\centering
\caption{Extrait représentatif du Product Backlog GerAI}
\label{tab:pbl}
\begin{tabularx}{\textwidth}{|C|l|L|l|C|}
\hline
\rowcolor{headerblue!25}\textbf{ID} & \textbf{Rôle} & \textbf{User Story} & \textbf{Priorité} & \textbf{Sprint} \\\hline
\rowcolor{rowgray}
US-01 & ADMIN\_RH & En tant qu'admin, je veux créer un employé avec code \texttt{EMP-NNNN} auto-incrémenté et provisioning Keycloak & Haute & 2 \\\hline
US-02 & EMPLOYE & En tant qu'employé, je veux déposer 12 types de congés avec vérification de solde & Haute & 3 \\\hline
\rowcolor{rowgray}
US-03 & CHEF & En tant que chef, je veux valider ou rejeter les demandes de mon équipe avec commentaire & Haute & 3 \\\hline
US-04 & EMPLOYE & En tant qu'employé, je veux connaître ma capacité d'endettement automatiquement & Haute & 4 \\\hline
\rowcolor{rowgray}
US-05 & COMITE & En tant que membre de la commission, je veux voter FAVORABLE/DÉFAVORABLE et suggérer un montant & Haute & 4 \\\hline
US-06 & EMPLOYE & En tant qu'employé, je veux soumettre des demandes par message à l'assistant IA & Moyenne & 5 \\\hline
\rowcolor{rowgray}
US-07 & CHEF & En tant que chef, je veux créer un projet, affecter des membres et voir l'avancement calculé & Haute & 6 \\\hline
US-08 & ADMIN\_RH & En tant qu'admin, je veux générer un rapport PDF/Excel de congés filtré par département & Haute & 7 \\\hline
\end{tabularx}
\end{table}

\subsection{Besoins non fonctionnels}

\begin{itemize}
  \item \textbf{Sécurité}~: chaque endpoint REST valide le JWT via
        \texttt{JwtAuthenticationConverter}~; la vérification des droits est
        déclarative via \texttt{@PreAuthorize} au niveau méthode.
  \item \textbf{Performance}~: les tableaux analytiques sont servis depuis le cache
        Spring Cache / Caffeine (TTL configurable)~; le frontend utilise la stratégie
        \texttt{ChangeDetectionStrategy.OnPush}.
  \item \textbf{Résilience}~: la défaillance d'un microservice n'affecte pas les
        autres~; les appels inter-services sont découplés par Kafka.
  \item \textbf{Reproductibilité}~: le schéma Oracle est versionné par 11 migrations
        Flyway~; \texttt{docker compose up} reconstruit l'intégralité de l'infrastructure.
  \item \textbf{Traçabilité}~: chaque transition de statut d'une demande est
        horodatée et associée à un valideur identifié par son \texttt{EMPLOYEE\_ID}.
\end{itemize}

% ─────────────────────────────────────────────────────────────
\section{Identification des acteurs}
% ─────────────────────────────────────────────────────────────

\begin{table}[H]
\centering
\caption{Acteurs du système et rôles Keycloak correspondants (realm \texttt{gerai})}
\label{tab:acteurs}
\begin{tabularx}{\textwidth}{|l|l|L|}
\hline
\rowcolor{headerblue!25}\textbf{Acteur} & \textbf{Rôle Keycloak} & \textbf{Responsabilités} \\\hline
\rowcolor{rowgray}
Administrateur RH & \texttt{ADMIN\_RH} & Cycle de vie employés, provisioning Keycloak, validation finale, rapports, attrition \\\hline
Chef de projet & \texttt{CHEF} & Avis hiérarchique, projets/tâches, recrutement, évaluations \\\hline
\rowcolor{rowgray}
Employé & \texttt{EMPLOYE} & Dépôt demandes, messagerie, assistant IA, profil \\\hline
Membre comité & \texttt{COMITE} & Vote sur les demandes de prêt soumises à la commission \\\hline
\end{tabularx}
\end{table}

% ─────────────────────────────────────────────────────────────
\section{Diagramme de cas d'utilisation global}
% ─────────────────────────────────────────────────────────────

\begin{figure}[H]
\centering
\begin{tikzpicture}[xscale=1.0, yscale=1.0]

  % ── Acteur Admin (gauche) ─────────────────────────────────
  \draw[thick] (-2.2,8.3) circle(0.3cm);
  \draw[thick] (-2.2,8.0)--(-2.2,7.2) coordinate(wA);
  \draw[thick] (wA)++(0,.1)--++(-.4,-.35) (wA)++(0,.1)--++(.4,-.35);
  \draw[thick] (wA)--++(-.28,-.55) (wA)--++(.28,-.55);
  \node[font=\small] at (-2.2,6.25) {Admin RH};

  % ── Acteur Chef ──────────────────────────────────────────
  \draw[thick] (-2.2,5.0) circle(0.3cm);
  \draw[thick] (-2.2,4.7)--(-2.2,3.9) coordinate(wC);
  \draw[thick] (wC)++(0,.1)--++(-.4,-.35) (wC)++(0,.1)--++(.4,-.35);
  \draw[thick] (wC)--++(-.28,-.55) (wC)--++(.28,-.55);
  \node[font=\small] at (-2.2,3.0) {Chef};

  % ── Acteur Employé ───────────────────────────────────────
  \draw[thick] (-2.2,1.8) circle(0.3cm);
  \draw[thick] (-2.2,1.5)--(-2.2,0.7) coordinate(wE);
  \draw[thick] (wE)++(0,.1)--++(-.4,-.35) (wE)++(0,.1)--++(.4,-.35);
  \draw[thick] (wE)--++(-.28,-.55) (wE)--++(.28,-.55);
  \node[font=\small] at (-2.2,-0.15) {Employé};

  % ── Acteur Comité (droite) ───────────────────────────────
  \draw[thick] (13,2.8) circle(0.3cm);
  \draw[thick] (13,2.5)--(13,1.7) coordinate(wK);
  \draw[thick] (wK)++(0,.1)--++(-.4,-.35) (wK)++(0,.1)--++(.4,-.35);
  \draw[thick] (wK)--++(-.28,-.55) (wK)--++(.28,-.55);
  \node[font=\small] at (13,0.7) {Comité};

  % ── Frontière système ────────────────────────────────────
  \begin{scope}[on background layer]
    \filldraw[draw=headerblue!60, fill=blue!3, rounded corners=6pt, dashed]
      (0,-0.5) rectangle (11.5,9.5);
    \node[font=\bfseries\small, headerblue] at (5.75,9.2) {Plateforme GerAI};
  \end{scope}

  % ── Cas d'utilisation ────────────────────────────────────
  % Module Employés
  \draw (4,8.8) ellipse(2.2cm and .45cm);
  \node[font=\scriptsize] at (4,8.8) {Gérer les employés};
  \draw (8.5,8.8) ellipse(2.2cm and .45cm);
  \node[font=\scriptsize] at (8.5,8.8) {Générer rapports / attrition};

  % Module Demandes
  \draw (4,7.3) ellipse(2.2cm and .45cm);
  \node[font=\scriptsize] at (4,7.3) {Valider les demandes};
  \draw (8.5,7.3) ellipse(2.2cm and .45cm);
  \node[font=\scriptsize] at (8.5,7.3) {Gérer projets et tâches};

  \draw (4,5.8) ellipse(2.2cm and .45cm);
  \node[font=\scriptsize] at (4,5.8) {Déposer une demande};
  \draw (8.5,5.8) ellipse(2.2cm and .45cm);
  \node[font=\scriptsize] at (8.5,5.8) {Recrutement / scoring CV};

  \draw (4,4.3) ellipse(2.2cm and .45cm);
  \node[font=\scriptsize] at (4,4.3) {Consulter suivi (timeline)};
  \draw (8.5,4.3) ellipse(2.2cm and .45cm);
  \node[font=\scriptsize] at (8.5,4.3) {Messagerie / assistant IA};

  \draw (6.25,2.5) ellipse(2.4cm and .45cm);
  \node[font=\scriptsize] at (6.25,2.5) {Voter sur une demande de prêt};

  \draw (6.25,1.0) ellipse(2.4cm and .45cm);
  \node[font=\scriptsize] at (6.25,1.0) {Consulter son profil};

  % ── Associations ─────────────────────────────────────────
  \draw[->] (-1.88,8.3) -- (1.8,8.8);
  \draw[->] (-1.88,8.1) -- (6.3,8.8);
  \draw[->] (-1.88,7.7) -- (1.8,7.3);
  \draw[->] (-1.88,4.7) -- (1.8,5.8);
  \draw[->] (-1.88,4.5) -- (1.8,4.3);
  \draw[->] (-1.88,4.2) -- (6.3,7.3);
  \draw[->] (-1.88,3.5) -- (6.3,4.3);
  \draw[->] (-1.88,1.8) -- (1.8,5.8);
  \draw[->] (-1.88,1.6) -- (1.8,4.3);
  \draw[->] (-1.88,1.4) -- (6.3,4.3);
  \draw[->] (-1.88,1.1) -- (6.3,1.0);
  \draw[->] (12.7,2.8) -- (8.65,2.5);

\end{tikzpicture}
\caption{Diagramme de cas d'utilisation global de la plateforme GerAI}
\label{fig:uc_global}
\end{figure}

% ─────────────────────────────────────────────────────────────
\section{Architecture microservices}
% ─────────────────────────────────────────────────────────────

\subsection{Vue d'ensemble}

\begin{figure}[H]
\centering
\begin{tikzpicture}[node distance=1.0cm]
  % Frontend
  \node[draw, rounded corners=5pt, fill=green!10, minimum width=11.5cm,
        minimum height=0.9cm, font=\small\bfseries] (ng) at (5.75,8.5)
    {Angular — composants \textit{standalone}, stratégie OnPush, \texttt{keycloak-angular}};

  % Services ligne 1
  \node[component] (emp)  at (0,   6.2) {employe-service\\:8081};
  \node[component] (dem)  at (2.9, 6.2) {demandes-service\\:8085};
  \node[component] (ch)   at (5.75,6.2) {chat-service\\:8086};
  \node[component] (not)  at (8.6, 6.2) {notification-service\\:8084};
  \node[component] (ana)  at (11.5,6.2) {analytics-service\\:8083};

  % Services ligne 2
  \node[component] (prj)  at (3.8, 4.6) {projets-service\\:8087};
  \node[component] (tch)  at (7.7, 4.6) {taches-service\\:8088};

  % Infrastructure
  \node[infra] (ora) at (0,   2.5) {Oracle 23ai\\:1521};
  \node[infra] (kc)  at (2.9, 2.5) {Keycloak 26\\:8080};
  \node[infra] (kaf) at (5.75,2.5) {Kafka\\:9092};
  \node[infra] (mh)  at (8.6, 2.5) {MailHog\\:1025};

  % Flèches frontend → services
  \foreach \t in {emp,dem,ch,not,ana,prj,tch}
    \draw[compArr] (ng.south) -- (\t.north);

  % Flèches services → BDD/Infra
  \foreach \t in {emp,dem,ch,ana,prj,tch}
    \draw[compArr] (\t.south) -- (ora.north);
  \draw[compArr] (emp.south) to[bend right=15] (kc.north);
  \draw[compArr] (dem.south) -- (kaf.north);
  \draw[compArr] (not.south) -- (kaf.north);
  \draw[compArr] (not.south) to[bend left=10] (mh.north);

  % WebSocket et OAuth2
  \draw[compArr,dashed,blue!60] (ng.east) to[bend right=25]
    node[midway,right,font=\scriptsize,blue!70]{WebSocket\\STOMP} (ch.north east);
  \draw[compArr,dashed,red!60] (ng.west) to[bend left=25]
    node[midway,left,font=\scriptsize,red!70]{OAuth2/OIDC} (kc.north west);
\end{tikzpicture}
\caption{Architecture microservices de GerAI~: services, protocoles et infrastructure}
\label{fig:archi}
\end{figure}

\begin{table}[H]
\centering
\caption{Microservices de GerAI~: ports, entités principales et responsabilités}
\label{tab:services}
\begin{tabularx}{\textwidth}{|l|C|L|}
\hline
\rowcolor{headerblue!25}\textbf{Service} & \textbf{Port} & \textbf{Entités principales et responsabilité} \\\hline
\rowcolor{rowgray}
\texttt{employe-service} & 8081 & \texttt{Employee}, \texttt{EmployeeDocument}~; cycle de vie, sync Oracle--Keycloak \\\hline
\texttt{demandes-service} & 8085 & \texttt{LeaveRequest}, \texttt{LoanRequest}, \texttt{DocumentRequest}, \texttt{TrainingRequest}, \texttt{AuthorizationRequest}, \texttt{ComiteVote}~; workflows RH \\\hline
\rowcolor{rowgray}
\texttt{chat-service} & 8086 & \texttt{Message}, \texttt{Conversation}, \texttt{ConversationParticipant}~; WebSocket STOMP, assistant Groq \\\hline
\texttt{notification-service} & 8084 & \texttt{Notification}~; consommateur Kafka, e-mail MailHog, push WebSocket \\\hline
\rowcolor{rowgray}
\texttt{analytics-service} & 8083 & \texttt{AttritionPredictionDTO}~; JasperReports (7 templates), scoring attrition \\\hline
\texttt{projets-service} & 8087 & \texttt{Project}, \texttt{Task}, \texttt{ProjectMember}, \texttt{Candidate}, \texttt{Interview}~; projets, tâches, recrutement \\\hline
\rowcolor{rowgray}
\texttt{taches-service} & 8088 & \texttt{Tache}~; gestion individuelle des tâches avec notifications \\\hline
\end{tabularx}
\end{table}

% ─────────────────────────────────────────────────────────────
\section{Schéma de base de données}
% ─────────────────────────────────────────────────────────────

La base Oracle 23ai comprend les tables suivantes, versionnées par 11 migrations Flyway (V1--V11)~:

\begin{table}[H]
\centering
\caption{Tables principales du schéma Oracle~: colonnes clés}
\label{tab:schema}
\begin{tabularx}{\textwidth}{|l|L|}
\hline
\rowcolor{headerblue!25}\textbf{Table} & \textbf{Colonnes significatives} \\\hline
\rowcolor{rowgray}
\texttt{EMPLOYEES} & \texttt{EMPLOYEE\_ID (PK)}, \texttt{USER\_ID (UNIQUE)} — UUID Keycloak, \texttt{EMPLOYEE\_CODE} — format \texttt{EMP-0042}, \texttt{DEPT\_ID}, \texttt{POSITION\_ID}, \texttt{MANAGER\_ID} (auto-jointure), \texttt{STATUS} (ACTIF/INACTIF/SUSPENDU) \\\hline
\texttt{LOAN\_REQUESTS} & \texttt{REQUEST\_ID (PK)}, \texttt{AMOUNT NUMBER(12,2)}, \texttt{DURATION\_MONTHS}, \texttt{STATUS} (EN\_ATTENTE/EN\_ETUDE/APPROUVE/REFUSE), \texttt{NEEDS\_COMMISSION NUMBER(1) DEFAULT 1} \\\hline
\rowcolor{rowgray}
\texttt{LEAVE\_REQUESTS} & \texttt{REQUEST\_ID (PK)}, \texttt{LEAVE\_TYPE\_ID}, \texttt{START\_DATE}, \texttt{END\_DATE}, \texttt{DAYS\_COUNT}, \texttt{STATUS}, \texttt{MANAGER\_COMMENT} \\\hline
\texttt{COMITE\_VOTES} & \texttt{VOTE\_ID (PK)}, \texttt{LOAN\_ID (FK)}, \texttt{MEMBER\_ID (FK)}, \texttt{VOTE} (FAVORABLE/DEFAVORABLE), \texttt{MONTANT\_SUGGERE}, \texttt{NB\_TRANCHES\_SUGGERES} \\\hline
\rowcolor{rowgray}
\texttt{PROJECTS} & \texttt{PROJECT\_ID (PK)}, \texttt{NAME}, \texttt{STATUS}, \texttt{PROGRESS\_PCT}, \texttt{CREATED\_BY (FK EMPLOYEES)} \\\hline
\texttt{V\_ALL\_DEMANDES} & Vue Oracle agrégeant les 5 tables de demandes~; utilisée par l'analytics et la requête d'attrition \\\hline
\end{tabularx}
\end{table}

% ─────────────────────────────────────────────────────────────
\section{Environnement de développement}
% ─────────────────────────────────────────────────────────────

\begin{table}[H]
\centering
\caption{Environnement matériel et logiciel}
\label{tab:env}
\begin{tabularx}{\textwidth}{|l|L|}
\hline
\rowcolor{headerblue!25}\textbf{Composant} & \textbf{Configuration} \\\hline
\rowcolor{rowgray}Processeur & Intel Core i7 11\textsuperscript{e} génération, 2,8 GHz \\\hline
RAM & 16 Go DDR4 \\\hline
\rowcolor{rowgray}Système & Windows 11 Professionnel 64 bits \\\hline
IDE backend & IntelliJ IDEA 2024 (Java 21, Spring Boot 3.x) \\\hline
\rowcolor{rowgray}IDE frontend & Visual Studio Code + Angular Extension Pack \\\hline
Test API & Postman \\\hline
\rowcolor{rowgray}Admin BDD & DBeaver (Oracle 23ai) \\\hline
Conteneurs & Docker Desktop~: \texttt{gvenzl/oracle-free:23-slim}, \texttt{quay.io/keycloak/keycloak:26.5.5}, \texttt{confluentinc/cp-kafka:7.5.3} \\\hline
\rowcolor{rowgray}Versioning & GitHub \\\hline
Rapport & Overleaf (compilation \LaTeX) \\\hline
\end{tabularx}
\end{table}

% ─────────────────────────────────────────────────────────────
\section{Choix technologiques justifiés}
% ─────────────────────────────────────────────────────────────

\subsection{Architecture~: microservices vs alternatives}

\begin{table}[H]
\centering
\caption{Comparaison des styles architecturaux}
\begin{tabularx}{\textwidth}{|l|C|C|C|}
\hline
\rowcolor{headerblue!25}\textbf{Critère} & \textbf{Monolithique} & \textbf{SOA/ESB} & \textbf{Microservices} \\\hline
\rowcolor{rowgray}Scalabilité indépendante & \ding{55} & Partielle & \checkmark \\\hline
Isolation des pannes & \ding{55} & Partielle & \checkmark \\\hline
\rowcolor{rowgray}Déploiement par domaine & \ding{55} & Partiel & \checkmark \\\hline
Cache différencié par service & Impossible & Difficile & \checkmark \\\hline
\rowcolor{rowgray}Complexité opérationnelle & \textbf{Faible} & Moyenne & Élevée \\\hline
\end{tabularx}
\end{table}

\noindent\textbf{Décision}~: l'\textit{analytics-service} (intensif en I/O Oracle)
doit pouvoir évoluer indépendamment du \textit{chat-service} (intensif en mémoire
WebSocket). Seule l'architecture microservices rend cette isolation possible sans
reconfiguration des autres services.

\subsection{Framework backend~: Spring Boot 3.x}

\begin{table}[H]
\centering
\caption{Comparaison des frameworks backend}
\begin{tabularx}{\textwidth}{|l|C|C|C|}
\hline
\rowcolor{headerblue!25}\textbf{Critère} & \textbf{Spring Boot 3.x} & \textbf{Quarkus} & \textbf{Micronaut} \\\hline
\rowcolor{rowgray}Intégration Keycloak native & \checkmark & \checkmark & Partielle \\\hline
Spring Kafka mature & \checkmark & Extension & Extension \\\hline
\rowcolor{rowgray}Spring Data JPA / Oracle & \checkmark & Hibernate ORM & Hibernate ORM \\\hline
Écosystème JasperReports & \checkmark & Limité & Limité \\\hline
\rowcolor{rowgray}Maturité Java 21 & \checkmark & \checkmark & \checkmark \\\hline
\end{tabularx}
\end{table}

\noindent\textbf{Décision}~: \texttt{spring-security-oauth2-resource-server}
configure la validation JWT Keycloak en quelques lignes YAML~;
\texttt{spring-kafka} fournit \texttt{@KafkaListener} et \texttt{KafkaTemplate}
éprouvés~; \texttt{spring-data-jpa} couvre la totalité des besoins Oracle avec
les requêtes natives utilisées dans le scoring d'attrition.

\subsection{Framework frontend~: Angular}

Angular est retenu pour sa cohérence architecturale indispensable à une
application de gestion complexe~: injection de dépendances intégrée, formulaires
réactifs (\texttt{ReactiveFormsModule}) avec validateurs conditionnels activés
dynamiquement selon le type de demande, routing déclaratif protégé par rôle
Keycloak, et stratégie \texttt{ChangeDetectionStrategy.OnPush} réduisant les
re-rendus dans les vues de listes paginées.

\subsection{Base de données~: Oracle 23ai Free}

Oracle est retenu pour ses fonctions analytiques SQL avancées (fenêtrage,
requêtes \texttt{WITH} récursives) utilisées dans la requête de scoring
d'attrition (\texttt{findAttritionRawData()}) et pour son adéquation avec
le parc technologique Oracle déjà maîtrisé chez Arabsoft. La version Free
déployée depuis \texttt{gvenzl/oracle-free:23-slim} élimine tout coût.

\subsection{Messagerie asynchrone~: Apache Kafka}

Kafka est préféré à RabbitMQ pour la durabilité des messages sur disque et la
possibilité de rejouer les événements~: si le \textit{notification-service} est
temporairement indisponible, les événements publiés sur \texttt{notification-events}
ne sont pas perdus. Les trois services producteurs (demandes, employe, projets)
publient des \texttt{NotificationEvent}~; le \textit{notification-service} les
consomme via \texttt{@KafkaListener} sur le groupe \texttt{notification-group}.

\subsection{Serveur d'identité~: Keycloak 26}

Keycloak est retenu face à Auth0 (coût SaaS) et WSO2 (complexité initiale) pour
son déploiement on-premise gratuit et son API d'administration REST complète~:
l'\textit{employe-service} crée les comptes Keycloak (\texttt{POST /admin/realms/gerai/users}),
affecte les rôles et réinitialise les mots de passe provisoires en trois appels
REST successifs.

% ─────────────────────────────────────────────────────────────
\section{Conclusion}
% ─────────────────────────────────────────────────────────────

Ce chapitre a formalisé les exigences de GerAI, décrit son architecture
microservices fondée sur sept services aux responsabilités délimitées, présenté
le schéma Oracle versionné par Flyway, et justifié chaque choix technologique
par comparaison argumentée d'alternatives. Ces décisions constituent le cadre
de référence de tous les développements décrits dans les chapitres suivants.
